\section{从 Internet-scale Video 学什么：Representation、Reward 与 Action}

MineCLIP、Voyager 和 Eureka 之后，课程转向更长期的研究方向。互联网视频便宜、动态、覆盖大量人类行为，却来自第三人称或不受控视角，通常没有 agent action、state 或 grounding label。问题不应被笼统称为“从视频模仿”，而要拆成三种监督目标：学 representation、学 reward、或推断 action 后做 behavior cloning。

\subsection{展望：两条 Foundation-Agent 路线}

讲者把未来方向分成两类：继续从 internet-scale video 学习，以及构建原生 multimodal foundation model。前者扩展行为先验，后者统一文本、图像、视频、声音与 action 接口；两者最终都要回到 active loop，接受 embodiment 和 environment feedback。

\lecturefigure{slide-43-looking-forward-title.jpg}{从当前系统过渡到 generalist-agent 长期方向}{官方录像恢复幻灯片 V043；00:37:36}

这页像一道章节门：前面的系统已经证明 reward、code 与 memory 可行，后面要追问怎样扩大数据和 modality。读图时不要把 chat/image/music 的成功直接外推到 agents；行动数据包含因果、控制和风险，不能只通过更多静态 token 自动获得。

\lecturefigure{slide-44-two-foundation-directions.jpg}{Internet-scale video learning 与 multimodal foundation models}{官方录像恢复幻灯片 V044；00:38:09}

\begin{importantbox}{两条路线最终在 action grounding 汇合}
Video learning 提供动态行为与人类策略，multimodal model 提供统一 token/interface；但系统仍需知道“哪个 action 由谁执行、对当前 embodiment 产生什么后果”。没有 active grounding，模型最多学到 observational correlation。
\end{importantbox}

\subsection{互联网视频的优势与缺失}

视频比文本更接近行为，因为它展示物体随时间变化、手如何接触工具、任务如何展开。可是观看者不知道相机外发生了什么，也看不到精确 action command、力、触觉或失败重试。\term{behavior cloning（行为克隆）}是用示范中的 observation--action pair 监督训练 policy；没有 action label 时，必须先恢复或近似这些 pair。

\lecturefigure{slide-45-internet-video-properties.jpg}{Internet video：便宜、动态、丰富，但缺少 action 与 grounding}{官方录像恢复幻灯片 V045；00:38:48}

\begin{knowledgebox}{术语消化：从视频学习的三条路径}
\begin{tabularx}{\textwidth}{@{}L{0.19\textwidth}X X@{}}
\toprule
路径 & 学到什么 & 关键缺口 \\
\midrule
Representation & 可迁移的时空与语义表征 & 表示好不等于会控制 \\
Reward learning & goal 与 trajectory 的相似度/价值 & reward bias 与 shortcut \\
Action / behavior cloning & 从 observation 预测可执行 action & 视频通常没有动作标签，且 embodiment 不同 \\
\bottomrule
\end{tabularx}
\end{knowledgebox}

\lecturefigure{slide-46-internet-video-example.jpg}{第三人称互联网视频展示行为，却隐藏控制信号}{官方录像恢复幻灯片 V046；00:39:39}

这个例子应先看手、对象和结果之间的时序，再问真正的 action 在哪里。像素只显示表面运动，不提供关节命令、施力、意图或相机外状态；同一视觉轨迹可能由多种控制产生。视频适合学 semantic dynamics，却需要额外假设才能变成某个 robot 的 motor supervision。

\teachervoice{讲者称互联网视频“cheap, dynamic, rich in human behavior”，随后马上补充被动观察者拿不到 action labels 与 grounding。这个转折很重要：数据规模解决 coverage，不自动解决 causal control；主动小猫问题在这里重新出现。}

\subsection{路径一：Video Representation Learning}

第一条路线先学习时空 representation，再用少量 domain-specific agent data fine-tune。模型可以通过 future prediction、masked modeling 或 video--text alignment 获得对象、动作阶段和场景变化的表示，但控制能力仍取决于下游 action data 与 interface。

\lecturefigure{slide-47-representation-from-video.jpg}{从视频预训练 representation，再适配 agent 任务}{官方录像恢复幻灯片 V047；00:40:18}

\begin{knowledgebox}{读图：预训练与控制微调承担不同职责}
图中的大规模视频阶段提供一般动态特征，后续少量 agent data 才把 representation 接到具体 action space。先比较冻结 encoder、fine-tune 与从零训练，再检查下游数据规模；若没有 action-grounded evaluation，只能说明 feature 有用，不能说明 policy 已学会任务。
\end{knowledgebox}

\subsection{路径二：从视频学习 Reward}

上一条路线只要求视频预训练出有用表示，本节进一步让视频直接参与任务优化：MineCLIP 把“轨迹看起来多像目标”变成 reward，指导 agent 在自己的环境里探索动作。它不要求视频与机器人共享精确 control，因此比逐动作模仿更容易跨接口；代价是必须验证视觉语义能否迁移，并防止 policy 主动利用 reward model 的盲点。

\lecturefigure{slide-48-reward-from-video.jpg}{Video-language 与 goal-conditioned representation 可形成 learned reward}{官方录像恢复幻灯片 V048；00:40:45}

\begin{knowledgebox}{读图：reward transfer 比 action transfer 更弱也更灵活}
视频只需告诉系统“结果看起来是否更接近目标”，具体 action 由 policy 在自己的 embodiment 中探索，因此可以跨控制接口；但 reward 只约束 outcome appearance 时，可能忽略路径、安全或隐含条件。读结果时必须同时检查 final state、trajectory 与 constraint。
\end{knowledgebox}

\subsection{路径三：恢复 Action 并做 Behavior Cloning}

第三条路线最接近 imitation learning：从视频推断 latent action 或使用模型标注 action，再训练 behavior policy。Video PreTraining（VPT）在 Minecraft 中通过少量带键鼠标签的数据训练 inverse dynamics model，再给大规模无标签视频补 action；这种方法依赖观察视角、动作空间和 embodiment 足够匹配。

\lecturefigure{slide-49-behavior-cloning-from-video.jpg}{从视频恢复 action label，再进行 behavior cloning}{官方录像恢复幻灯片 V049；00:42:33}

若 $z_t$ 表示从相邻帧推断的 latent action，真实 agent action 为 $a_t$，可以写成：
\[
z_t\sim q_\psi(z\mid o_t,o_{t+1}),\qquad
a_t\sim p_\omega(a\mid z_t,e),\qquad
\mathcal{L}_{\mathrm{BC}}=-\sum_t\log\pi_\theta(a_t\mid o_{\le t},g).
\]
其中 $q_\psi$ 是 inverse/latent-action model，$p_\omega$ 把 latent action 映射到 embodiment $e$ 的可执行动作，$\pi_\theta$ 是 behavior policy，$g$ 是目标。若视频与 agent 的 embodiment、camera 或 action support 不匹配，$p_\omega$ 会成为主要误差源。

\begin{warningbox}{Observation similarity 不等于 action identifiability}
仅凭两帧图像通常不能唯一恢复施力、速度和隐藏接触；第三人称视频还存在相机运动与遮挡。Action reconstruction 应报告 uncertainty，并允许 downstream policy 通过环境交互修正，而不是把伪标签当作 ground truth。
\end{warningbox}

\teachervoice{讲者把 representation、reward 与 action learning 明确分成三条路线，是为了避免“用视频训练 agent”成为空泛口号。三者对 action label 的需求逐步增加：representation 最弱，reward 居中，behavior cloning 最强；系统应按可获得监督选择接口。}

\subsection{本章小结}

Internet video 提供廉价、动态、覆盖广的人类行为数据，却缺少 action label 与 grounding。Representation learning 学特征，reward learning 学目标相似度，behavior cloning 尝试恢复动作；三条路线的监督强度、迁移假设和失败模式不同。规模可以扩大先验，但 active environment feedback 仍是验证行动因果的关键。

